\documentclass[
         english, % thesis will be in english
    % signatureplaces, % adds signature places in the title page
    monochrome, % makes all link colors black
]{VUMIFSEMasterThesis}
\usepackage{float}
\usepackage{wrapfig2}
\usepackage{hyperref}
\usepackage{algorithmicx}
\usepackage{algorithm}
\usepackage{algpseudocode}
\usepackage{amsfonts}
\usepackage{amsmath}
\usepackage{bm}
\usepackage{caption}
\usepackage{color}
\usepackage{graphicx}
\usepackage{listings}
\usepackage{subcaption}
\usepackage{biblatex}


% Title page
\university{VILNIUS UNIVERSITY}
\faculty{FACULTY OF MATHEMATICS AND INFORMATICS}
\department{SOFTWARE ENGINEERING MASTERS STUDY PROGRAM}
\papertype{Master thesis}
\title{End-to-End Neural Text-to-Speech Synthesis for the Lithuanian Language}
\titleineng{}
\author{Vytautas Lėveris}
% \secondauthor{Vardonis Pavardonis}   % Add second author
\supervisor{Dr. Gražina Korvel}
\reviewer{...} % Missing
\date{Vilnius – \the\year}

\bibliography{bibliography}

\begin{document}
\maketitle

%% Acknowledgements of individuals and/or institutions.
% \sectionnonumnocontent{}
% \vspace{7cm}
% \begin{center}
%     Padėkos asmenims ir/ar organizacijoms
% \end{center}

%\begin{lithuanian}
%\sectionnonumnocontent{Santrauka}
%Santraukose lietuvių ir anglų kalbomis glaustai aprašomas darbo turinys:
%pristatoma nagrinėta problema ir padarytos išvados. Santraukų gale nurodomi
%darbo raktiniai žodžiai. Santraukos lietuvių ir anglų kalbomis rašomos
%atskiruose puslapiuose. Kiekvienos jų apimtis -- ne daugiau kaip 0,5 puslapio.
%Automatiškai naudojamos lietuviškos kabutės: \enquote{tekstas}.

% Specify up to 5 most important topic keywords.
% A keyword can have multiple words.
%\raktiniaizodziai{raktinis žodis 1, raktinis žodis 2, raktinis žodis 3, raktinis žodis 4, raktinis žodis 5}
%\end{lithuanian}

%\begin{english}
%\sectionnonumnocontent{Summary}
%English summary. English quotes are used automatically: \enquote{tekstas}.

%\keywords{keyword 1, keyword 2, keyword 3, keyword 4, keyword 5}
%\end{english}

\tableofcontents

\sectionnonum{Introduction}

\subsection*{Motivation}

Text-to-Speech (TTS) synthesis is a field of speech technology that focuses on converting written text into natural-sounding human speech. Over the past decades, TTS systems have evolved from simple rule-based mechanisms producing robotic and monotonous speech into highly advanced neural systems capable of generating expressive, fluent, and human-like audio. Today, TTS technology plays an important role in modern human–computer interaction and has become an integral component of many everyday applications, including virtual assistants (e.g. Siri or Alexa), accessibility tools, navigation systems, audiobooks, customer service automation, and language learning or text translation platforms~\cite{kavitha2024enhancing}.

The growing demand for natural speech interfaces is driven by rapid advancements in artificial intelligence (AI), particularly in deep learning and neural network architectures. Neural TTS systems have significantly improved speech naturalness, prosody modeling, and pronunciation accuracy compared to earlier approaches. These systems learn the mapping from text directly to speech waveforms or intermediate acoustic representations, reducing the need for handcrafted rules and linguistic expertise. Modern TTS solutions are more scalable and adaptable across different languages and speaking styles.

Despite the impressive progress achieved for widely spoken languages such as English, Chinees or Spanish, many smaller or low-resource languages remain underrepresented in TTS research and development. Lithuanian, a Baltic language with complex morphology, rich inflection, and distinctive phonetic characteristics, is one such language. Developing high-quality neural TTS systems for Lithuanian presents both linguistic and technical challenges, making it an interesting and relevant research topic.

This Master’s thesis focuses on end-to-end neural text-to-speech synthesis for the Lithuanian language, with the goal of evaluating state-of-the-art TTS models, training them on Lithuanian speech data, and experimentally assessing their performance and suitability for this language.

\subsection*{Applications of Text-to-Speech Technology}

The importance of TTS technology extends beyond convenience and entertainment. One of its most impactful application areas is accessibility. For individuals with visual impairments, reading difficulties, or certain motor disabilities, TTS systems enable access to digital content, educational materials, and communication tools~\cite{chalamandaris2010unit, reddy2023speech}. Screen readers and voice-based interfaces rely heavily on high-quality speech synthesis to provide an inclusive user experience.

Another major application domain is human–computer interaction. Voice assistants such as smart speakers, mobile assistants, and in-car systems use TTS to communicate information naturally and efficiently~\cite{mohasi2006text}. As voice-based interaction becomes more prevalent, users increasingly expect synthesized speech to be expressive, emotionally appropriate, and indistinguishable from human speech.

TTS is also widely used in education and language learning, where it can provide pronunciation examples, reading assistance, and interactive learning experiences. In media and content creation, TTS enables the automated generation of audiobooks, news narration, podcasts, and video voiceovers, significantly reducing production costs and time.

In addition, TTS technology has applications in telecommunications, customer support automation, assistive robotics, and game development, where dynamic and context-aware speech synthesis enhances realism and user engagement. These diverse applications highlight the societal and economic value of advancing TTS research, particularly for languages that currently lack high-quality solutions.

\subsection*{Challenges and Motivation for Lithuanian TTS}

Lithuanian presents specific challenges for TTS systems due to its rich morphology, complex stress patterns, and phonetic distinctions. The language exhibits a large number of inflected forms, which increases vocabulary size and complicates text normalization and pronunciation modeling. Stress placement in Lithuanian can vary and plays an important role in meaning and naturalness of speech.

Lithuanian is considered a low-resource language in the context of speech technology, as publicly available, high-quality speech datasets are limited compared to major world languages. This makes it particularly interesting to study how modern neural TTS architectures perform under data constraints and what strategies can be used to improve synthesis quality.

Developing effective TTS systems for Lithuanian is not only academically relevant but also socially meaningful. High-quality Lithuanian TTS could support accessibility, education, digital preservation of the language, and the development of local AI-driven products and services.

\subsection*{Research Object}

investigate and evaluate end-to-end neural text-to-speech synthesis models for the Lithuanian language.

\subsection*{Research Goal}

Experiment and understand neural TTS performance for morphologically rich, low-resource languages and to provide insights that may support future development of Lithuanian speech synthesis systems.

\subsection*{Tasks to achieve the goal}
\begin{enumerate}
\item a review of classical and neural TTS approaches.
\item An analysis of modern end-to-end TTS models.
\item Preparation and processing of Lithuanian language datasets.
\item Training and experimentation with selected TTS models.
\item Evaluation of synthesized speech quality using objective and/or subjective metrics.
\end{enumerate}

\subsection*{Research Methods}
\begin{itemize}
\item ;
\item ;
\item ;
\item ;
\item .
\end{itemize}


\section{Overview of Text-to-Speech Synthesis}

This chapter reviews the development of TTS synthesis from rule-based and signal-processing methods to modern neural approaches. Section~\ref{sec:tts-overview} introduces TTS as a problem, defines its core components, and outlines the key challenges that persist regardless of the underlying technology. Sections~\ref{sec:classical} and~\ref{sec:neural} then trace the evolution of synthesis methods chronologically, covering classical approaches --- formant, concatenative, and statistical parametric synthesis --- followed by the neural architectures that have since become dominant.

\subsection{Text-to-Speech Synthesis as a Problem}
\label{sec:tts-overview}

Text-to-speech synthesis is the task of converting an arbitrary sequence of written characters into intelligible and natural-sounding speech. At a high level, a TTS system must solve two coupled sub-problems: determining \emph{what} to say (linguistic interpretation of the text) and \emph{how} to say it (acoustic realisation). The first sub-problem involves text normalisation, lexical disambiguation, grapheme-to-phoneme conversion, and prosody prediction; the second involves generating a waveform that encodes the target phonemes with appropriate pitch, timing, and voice quality.

A fundamental difficulty is that the mapping from text to speech is \emph{one-to-many}: a single written sentence can be pronounced in infinitely many valid ways differing in pitch range, speaking rate, emphasis, and emotional colouring, none of which is encoded in the orthographic input. Human speakers resolve this ambiguity through pragmatic context, intention, and personality --- information that is absent from raw text. TTS systems must therefore either predict a single plausible realisation from a learned prior, or explicitly model the space of possible prosodic realisations, a problem that remains an active area of research~\cite{hasanabadi2023overview}.

Additional challenges arise from the diversity of natural language. Text normalization must handle numerals, abbreviations, acronyms, foreign words, and domain-specific terminology, all of which require language- and context-dependent expansion rules. Grapheme-to-phoneme conversion is non-trivial in languages with irregular spelling, while languages with lexical stress or tone must correctly predict these features to avoid changing word meaning. Sentence-level prosody --- the placement of emphasis, phrasing boundaries, and intonational contours --- depends on syntactic structure, information structure, and discourse context, and predicting it reliably from text alone remains an open problem.

Computational cost is another persistent concern. High-quality neural waveform generation is computationally intensive: a model that produces 24~kHz audio must effectively predict 24,000 samples per second, each dependent on prior context. Autoregressive models address this by generating samples sequentially, at the cost of latency; parallel models can synthesise in a single forward pass but require more complex training procedures. Striking the right balance between quality, latency, and hardware requirements is particularly important for deployment on edge devices or in real-time interactive applications~\cite{kong2020hifi}.

Challenges are further amplified for low-resource languages. Most advances in TTS have been driven by research on English and a small number of other well-resourced languages, for which hundreds of hours of studio-quality recordings and large pronunciation lexicons are readily available. For low-resource languages such as Lithuanian, publicly available speech corpora are scarce and often limited in speaker diversity, recording conditions, and stylistic variety. Morphologically rich languages present an additional difficulty: the large number of inflected word forms increases lexical sparsity, complicates text normalisation, and reduces unit coverage in corpus-based approaches. As a result, models trained on low-resource data tend to produce less natural speech and generalise poorly to out-of-vocabulary words~\cite{hasanabadi2023overview}.

\subsection{Classical Approaches to TTS Synthesis}
\label{sec:classical}

A complete TTS system consists of two main processing stages. The \emph{front end} (text analysis) converts raw input text into a linguistic representation by performing text normalization, part-of-speech tagging, grapheme-to-phoneme conversion, and prosody prediction. The \emph{back end} (speech synthesis) uses this representation to generate an audio waveform. Classical TTS systems differ primarily in how the back end operates.

\subsubsection{Formant Synthesis}

Formant synthesis is one of the earliest approaches to machine speech production. It models the human vocal tract as a series of acoustic resonances known as formants — the first three of which (F1, F2, F3) largely determine vowel quality. A formant synthesizer generates speech by exciting a bank of resonant filters with either a periodic glottal source signal for voiced sounds, or a noise source for fricatives and other unvoiced sounds, and by varying the filter parameters over time according to the phonemic transcription and prosody specification provided by the front end.

The most influential formant synthesizer was described by Klatt~\cite{klatt1980software}, who presented a cascade/parallel architecture in which voiced sounds pass through resonators arranged in cascade (modelling the normal vocal tract transfer function), while nasals and some fricatives are handled by a parallel filter bank. The synthesizer was controlled by approximately 40 time-varying parameters — including fundamental frequency, voicing amplitude, formant frequencies and bandwidths for F1--F5, and source spectral tilt — updated at a frame rate of 5~ms. Commercial products such as DECTalk were built on this foundation and found widespread use in assistive technology throughout the 1980s and 1990s.

Formant synthesis offers practical advantages: the synthesizer itself is compact, requires no recorded speech database, and allows precise control over every acoustic parameter. However, producing natural-sounding output requires extensive phonological and phonetic expertise, and the handcrafted rules must be re-developed from scratch for each new language. In practice, formant-based systems produce speech that sounds robotic and monotonous, and they have largely been superseded by data-driven approaches~\cite{reddy2023speech}.

\subsubsection{Concatenative Synthesis}

Concatenative synthesis constructs speech by selecting and joining pre-recorded speech segments stored in a database, relying on the naturalness inherent in real human recordings rather than acoustic modelling.

The simplest concatenative approach uses \emph{diphones} as the basic unit~\cite{dutoit1997high}. A diphone spans from the middle of one phoneme to the middle of the following one, thereby capturing the coarticulatory transition between them. Because the inventory of distinct phoneme pairs in a language is relatively small — on the order of a few thousand — a complete diphone set can be recorded in a short session with a single speaker. After concatenation, prosodic modification is applied to match the target pitch contour and duration. The most widely used technique for this purpose is TD-PSOLA (Time-Domain Pitch Synchronous Overlap and Add)~\cite{moulines1990pitch}, which time-scales and pitch-shifts individual pitch periods while preserving the overall spectral envelope.

A more capable variant is \emph{unit selection} synthesis~\cite{hunt1996unit}. Rather than a minimal diphone inventory, unit selection systems index a large corpus of natural speech — typically several hours of recordings — and store every available acoustic unit (phones, half-phones, diphones, syllables, or words). At synthesis time, a Viterbi search selects the sequence of units that minimises a combined cost:
\begin{equation}
C = \sum_{i} \bigl[ C_t(u_i,\, t_i) + C_j(u_i,\, u_{i-1}) \bigr],
\end{equation}
where $C_t$ is the \emph{target cost} measuring how well unit $u_i$ matches the desired phonetic and prosodic target $t_i$, and $C_j$ is the \emph{join cost} measuring the spectral discontinuity at the boundary between consecutive units~\cite{hunt1996unit}. When a closely matching unit is found, no prosodic modification is required, and the resulting speech can be remarkably natural. Quality degrades, however, when required units are absent from the database or joins are acoustically mismatched.

Concatenative synthesis dominated commercial TTS production throughout the 1990s and 2000s and remains in use today in contexts where reliability and low latency are priorities. Because synthesis amounts to selecting and concatenating pre-recorded segments, no computationally intensive model inference is required: once the Viterbi search completes, waveform assembly is trivial and can run in real time on modest hardware without a GPU. This makes concatenative systems attractive for embedded or on-device deployment, interactive voice response (IVR) systems, and other latency-sensitive applications~\cite{chalamandaris2010unit}.

Despite these practical strengths, concatenative synthesis has several notable disadvantages:
\begin{itemize}
  \item \textbf{Database size.} Unit selection systems require large, carefully recorded speech corpora — often several hours of studio-quality audio from a single speaker — resulting in databases that can reach tens of gigabytes and are expensive to produce and maintain.
  \item \textbf{Fixed voice.} The system is tied to the voice of the recorded speaker; changing the speaker requires recording and indexing an entirely new corpus.
  \item \textbf{Limited prosodic flexibility.} Expressive or stylistically varied speech is difficult to achieve, since prosodic modification beyond the range of the recorded material introduces perceptible artefacts.
  \item \textbf{Coverage gaps.} For morphologically rich languages such as Lithuanian, the large number of distinct inflected word forms increases the probability that certain phonetic contexts are absent from even a large corpus, forcing the system to rely on poorly matching units or acoustically disruptive concatenations~\cite{chalamandaris2010unit}.
  \item \textbf{Degraded quality for out-of-coverage input.} When required units are missing or joins are spectrally discontinuous, output quality can drop sharply and unpredictably, making robustness difficult to guarantee across unrestricted input text.
\end{itemize}

\subsubsection{Statistical Parametric Speech Synthesis}

Statistical parametric speech synthesis (SPSS) addresses the database size and coverage limitations of concatenative methods by replacing the unit inventory with a compact statistical model trained on speech data. The dominant approach within this paradigm uses Hidden Markov Models (HMMs) and is embodied in the HTS (HMM-based Speech Synthesis) system~\cite{zen2009statistical, tokuda2013speech}.

In HMM-based TTS, each context-dependent phoneme is modelled by a left-to-right HMM whose states emit probability distributions over acoustic feature vectors. The features are extracted frame-by-frame from training recordings using a vocoder (e.g. STRAIGHT or WORLD) and typically include mel-generalized cepstral (MGC) coefficients representing the spectral envelope, log fundamental frequency F0, and aperiodicity measures. Because the space of possible phonetic contexts — accounting for surrounding phonemes, syllable position, stress, and sentence type — vastly exceeds the amount of training data, contexts are clustered using decision trees, allowing model parameters to be shared across similar contexts and alleviating data sparsity~\cite{zen2009statistical}.

At synthesis time, the sequence of context-dependent HMMs for the input utterance is determined by the front end, and a maximum-likelihood parameter generation (MLPG) algorithm~\cite{tokuda2013speech} is applied to produce a smooth trajectory of acoustic features that jointly satisfies constraints on the static features and their delta and delta-delta coefficients. The resulting feature sequence is then passed to the vocoder to reconstruct the waveform.

A notable strength of HMM-based synthesis is the range of techniques available for controlling and adapting the output voice, which became one of the main practical advantages of the paradigm~\cite{tokuda2013speech}.

\textbf{Speaker adaptation.} Starting from an average voice model (AVM) trained on data from many speakers, the model can be adapted to a target speaker using only a small amount of adaptation data. Maximum likelihood linear regression (MLLR) estimates a set of linear transformations applied to the HMM means (and optionally variances) so that the likelihood of the adaptation utterances is maximised~\cite{zen2009statistical}. A constrained variant, CMLLR (feature-space MLLR), applies a single affine transformation to the observation space rather than to individual model parameters, which is more robust when adaptation data is very limited. Maximum a posteriori (MAP) adaptation is an alternative that interpolates the adapted parameters with the AVM prior, preventing overfitting to sparse adaptation sets.

\textbf{Eigenvoice adaptation.} Eigenvoice methods~\cite{tokuda2013speech} construct a low-dimensional speaker space from the principal components of a set of training speaker models. A new speaker is represented as a point in this space, and adaptation reduces to estimating a small weight vector rather than a full set of model parameters. This makes it possible to adapt to a new voice from very few utterances --- sometimes a single sentence --- by searching for the speaker weights that maximise the likelihood of the adaptation data.

\textbf{Model interpolation.} Speaking style, emotion, or voice characteristics can be controlled by linearly interpolating between two or more separately trained HMM sets. For example, blending a neutral-style model with an expressive-style model at varying weights produces a smooth continuum of output styles~\cite{zen2009statistical}. The same mechanism can interpolate between speaker models to generate voices with intermediate characteristics, enabling synthetic voices that are not tied to any single recorded speaker.

Compared to concatenative synthesis, SPSS offers a markedly different set of trade-offs. On the practical side, a trained HMM system is orders of magnitude more compact than a unit selection database: the statistical model typically occupies a few megabytes, whereas a unit selection corpus can reach tens of gigabytes. Coverage is no longer a hard constraint --- the model can generalise to phonetic contexts not seen during training through decision-tree clustering and parameter sharing, which is particularly valuable for morphologically rich languages with large inflected vocabularies. Changing the output voice or speaking style does not require recording a new corpus; instead, it can be achieved through the adaptation and interpolation techniques described above. On the other hand, the naturalness ceiling of HMM-based synthesis is lower than that of a well-matched unit selection system: when concatenative synthesis finds closely matching units, the resulting speech is drawn directly from real human recordings and can sound highly natural, whereas HMM output is always the product of statistical averaging and vocoder reconstruction. In practice, HMM-based synthesis was therefore preferred when flexibility, low storage, and multi-speaker capability were priorities, while concatenative synthesis remained competitive for high-quality single-voice applications with controlled vocabulary~\cite{zen2009statistical, tokuda2013speech}.

The main drawback of HMM-based synthesis is that maximum likelihood estimation inherently averages over the training data, producing over-smoothed parameter trajectories that result in muffled speech lacking fine spectral detail. Additional artefacts arise from the vocoder-based waveform reconstruction, which introduces a characteristic buzziness. These shortcomings were the primary motivation for the shift toward neural network-based approaches~\cite{hasanabadi2023overview}.

\subsection{Emergence of Neural TTS Systems}
\label{sec:neural}

The limitations of statistical parametric synthesis --- over-smoothed trajectories and vocoder artefacts --- motivated a shift toward deep learning approaches that learn richer acoustic representations directly from data. This transition unfolded in several stages: first replacing individual SPSS components with neural networks, then developing fully neural vocoders, and finally building systems that map text to waveforms within a single model.

\subsubsection{Early DNN-Based Acoustic Models}

The first wave of neural TTS research retained the overall SPSS pipeline structure but replaced the HMM output distributions with deep neural networks (DNNs)~\cite{zen2013statistical}. A feed-forward DNN maps a linguistic feature vector --- encoding phoneme identity, position, stress, and sentential context --- directly to acoustic feature vectors, bypassing the need for decision-tree clustering. Training by minimising mean squared error over acoustic parameters produced smoother predictions than HMMs, but frame-independent DNN predictions still averaged over the training distribution and therefore suffered from the same over-smoothing problem.

Subsequent work introduced recurrent architectures, particularly Long Short-Term Memory (LSTM) networks, to model temporal dependencies across frames. RNN-based acoustic models produced more coherent parameter trajectories and established sequence-to-sequence modelling as the central paradigm for neural TTS~\cite{zen2013statistical}.

\subsubsection{Neural Vocoders}

Parallel to the development of neural acoustic models, a separate line of research addressed waveform reconstruction. Traditional vocoders (STRAIGHT, WORLD) introduce characteristic buzziness through their source--filter decomposition; neural vocoders instead learn to generate waveforms directly from an acoustic conditioning signal.

\paragraph{Autoregressive vocoders.}
WaveNet~\cite{oord2016wavenet}, proposed by van den Oord et al., was the first neural vocoder to match the naturalness of real recordings in listening tests. It is an autoregressive model that generates each audio sample conditioned on all previous samples and on a local acoustic conditioning signal, using stacks of dilated causal convolutions to maintain a large receptive field without excessive depth. Despite its quality, WaveNet is computationally expensive: synthesising one second of 16~kHz audio requires 16,000 sequential forward passes through the network.

WaveRNN~\cite{kalchbrenner2018efficient} reduced this cost by replacing the convolutional stack with a single gated recurrent unit operating on quantised samples, combined with subscale prediction and sparsification techniques, achieving real-time synthesis on mobile hardware.

\paragraph{Non-autoregressive vocoders.}
To eliminate sequential generation entirely, flow-based and generative adversarial network (GAN) architectures were developed. WaveGlow~\cite{prenger2019waveglow} uses normalising flows --- a sequence of invertible transformations with tractable Jacobians --- to transform a Gaussian noise vector into a waveform in a single forward pass conditioned on a mel-spectrogram, trained by directly maximising the data log-likelihood.

GAN-based vocoders pushed quality and speed further. MelGAN~\cite{kumar2019melgan} introduced a fully convolutional generator conditioned on mel-spectrograms and trained adversarially against multi-scale discriminators operating at different temporal resolutions. HiFi-GAN~\cite{kong2020hifi} extended this design with multi-period discriminators that explicitly model periodic structures in speech, and a multi-receptive-field fusion generator that processes waveforms at multiple pattern lengths simultaneously. HiFi-GAN produces audio quality comparable to WaveNet while running orders of magnitude faster than real time, and has become the de facto standard vocoder in modern TTS pipelines.

\subsubsection{Two-Step TTS Pipelines}

Combining a neural acoustic model (predicting mel-spectrograms from text) with a neural vocoder (reconstructing waveforms from mel-spectrograms) defines the dominant two-step TTS architecture. The two components are trained independently, which gives flexibility --- either part can be improved or replaced without retraining the other --- but also means the acoustic model is optimised for spectrogram reconstruction loss rather than perceptual waveform quality.

\paragraph{Autoregressive acoustic models.}
Tacotron~\cite{wang2017tacotron} was the first sequence-to-sequence acoustic model to synthesise intelligible speech directly from characters. It uses an encoder--attention--decoder architecture in which a content-based attention mechanism learns the monotonic alignment between input characters and output spectrogram frames, eliminating the need for hand-crafted phoneme dictionaries or explicit duration models. The decoder predicts spectrogram frames autoregressively, one or more at each step.

Tacotron~2~\cite{shen2018natural} refined this design with a more stable decoder and conditioned WaveNet directly on predicted mel-spectrograms, achieving a mean opinion score (MOS) close to that of natural speech. The autoregressive decoder nevertheless generates frames sequentially, making synthesis latency proportional to utterance length and making the model sensitive to alignment failures on longer or unusual inputs.

\paragraph{Non-autoregressive acoustic models.}
FastSpeech~\cite{ren2019fastspeech} addressed the latency and robustness limitations of autoregressive decoders by introducing a feed-forward Transformer that generates all output frames in parallel. A duration predictor, trained with durations extracted from a pre-trained Tacotron teacher model, determines how many frames each input phoneme should expand to. The phoneme representations are expanded accordingly and processed by a parallel Transformer decoder. FastSpeech achieved approximately 38-times faster spectrogram generation than Tacotron~2 with comparable perceptual quality.

FastSpeech~2~\cite{ren2021fastspeech} removed the dependence on a teacher model by predicting pitch, energy, and duration directly from the phoneme encoder output using dedicated variance predictors, trained with ground-truth values extracted from the training data. The explicit prosodic variables also enable controllable synthesis, allowing independent adjustment of speaking rate, pitch range, and energy at inference time.

\subsubsection{End-to-End Models}

Training the acoustic model and vocoder independently introduces a train--test mismatch: the vocoder is conditioned on natural mel-spectrograms during training but on predicted, imperfect ones at inference. End-to-end models eliminate this mismatch by jointly optimising the full pipeline from text to waveform.

VITS (Variational Inference with adversarial learning for Text-to-Speech)~\cite{kim2021conditional} is the most widely adopted single-stage model. It combines a variational autoencoder (VAE), a normalising flow network, and adversarial training within a unified framework. A posterior encoder maps a mel-spectrogram to a latent representation; a flow-based prior network learns the conditional distribution of latent variables given the text; and a HiFi-GAN decoder generates the waveform from the sampled latent. The adversarial loss and reconstruction loss are optimised jointly, so the decoder is trained on its own outputs from the start and the prior network is pushed to match the true posterior. At inference, the prior network samples a latent variable conditioned only on the input text and the decoder synthesises the waveform directly, with no intermediate spectrogram. VITS achieved higher MOS scores than Tacotron~2 with WaveGlow in subjective evaluations while producing speech in a single forward pass.

%\sectionnonum{Rezultatai}
%Rezultatų skyriuje išdėstomi pagrindiniai darbo rezultatai: kažkas išanalizuota,
%kažkas sukurta, kažkas įdiegta. Tarpinių žingsnių išdavos skirtos užtikrinti galutinio
%rezultato kokybę neturi būti pateikiami šiame skyriuje. Kalbant informatikos termi-
%nais, šiame skyriuje pateikiama darbo išvestis, kuri gali būti įvestimi kituose panašios
%tematikos darbuose. Rezultatai pateikiami sunumeruotų (gali būti hierarchiniai) sąrašų
%pavidalu. Darbo rezultatai turi atitikti darbo tikslą.

%\sectionnonum{Išvados}
%\begin{enumerate}[labelindent=0pt]
%    \item Išvadų skyriuje daromi nagrinėtų problemų sprendimo metodų palyginimai, siūlomos
%rekomendacijos, akcentuojamos naujovės.
%    \item Išvados pateikiamos sunumeruoto (gali būti hierarchinis) sąrašo pavidalu.
%    \item Darbo išvados turi atitikti darbo tikslą.
%\end{enumerate}

\printbibliography[heading=bibintoc]  % Šaltinių sąraše nurodoma panaudota
% literatūra, kitokie šaltiniai. Abėcėlės tvarka išdėstomi darbe panaudotų
% (cituotų, perfrazuotų ar bent paminėtų) mokslo leidinių, kitokių publikacijų
% bibliografiniai aprašai. Šaltinių sąrašas spausdinamas iš naujo puslapio.
% Aprašai pateikiami netransliteruoti. Šaltinių sąraše negali būti tokių
% šaltinių, kurie nebuvo paminėti tekste (LaTeX tai sutvarko automatiškai).
% Šaltinių sąraše rekomenduojame necituoti savo kursinio darbo, nes tai nėra
% oficialus literatūros šaltinis. Jei tokių nuorodų reikia, pateikti jas tekste.

% \sectionnonum{Sąvokų apibrėžimai}
%\sectionnonum{Santrumpos}
%Sąvokų apibrėžimai ir santrumpų sąrašas sudaromas tada, kai darbo tekste
%vartojami specialūs paaiškinimo reikalaujantys terminai ir rečiau sutinkamos
%santrumpos.

% Appendices
% Prieduose gali būti pateikiama pagalbinė, ypač darbo autoriaus savarankiškai
% parengta, medžiaga. Savarankiški priedai gali būti pateikiami ir
% kompaktiniame diske. Priedai taip pat numeruojami ir vadinami. Darbo tekstas
% su priedais susiejamas nuorodomis.
%\appendix{Neuroninio tinklo struktūra}

%\begin{figure}[H]
%    \centering
%    \includegraphics[scale=0.5]{img/MLP}
%    \caption{Paveikslėlio pavyzdys}
%    \label{img:mlp}
%\end{figure}


%\appendix{Eksperimentinio palyginimo rezultatai}

% tablesgenerator.com - converts calculators (e.g. excel) tables to LaTeX
%\begin{table}[H]\footnotesize
%  \centering
%  \caption{Lentelės pavyzdys}
%  {\begin{tabular}{|l|c|c|} \hline
%    Algoritmas & $\bar{x}$ & $\sigma^{2}$ \\
%    \hline
%    Algoritmas A  & 1.6335    & 0.5584       \\
%    Algoritmas B  & 1.7395    & 0.5647       \\
%    \hline
%  \end{tabular}}
%  \label{tab:table example}
%\end{table}

\end{document}
